\documentclass{jsarticle}
\usepackage{amsmath}
\usepackage{amssymb}
\begin{document}
\title{}
\author{}
\date{}
\maketitle

宇宙と異次元空間において、原子間距離から、微細構造定数により、逆関数には、
ゼータ関数がこの逆関数において、原子から宇宙空間においてのノルムが
ホログラム空間で、接平面の曲平面の２次元射影空間のミンコフスキー空間
がこのノルム空間の機構を表している。最小単位が逆関数により、無限空間を生み出す。
この無限空間は、指数定理により、相加相乗平均の対数関数による、指数関数からの
双曲点から、臨界値により勾配率が０の等加速度と等速度運動においての、
双曲微分から、この接空間の曲率を求めると、開集合により、
電子によるオーヴィタル雲から異次元空間がつくり出されている。
 
  原子核における中間子論は、この原子とオーヴィタル雲を形成している、
電子雲をゼータ関数と量子群を加群とする機構自体が、この中間子を形成している。
原子関数のオーヴィタル雲が自身にエントロピー値を保有している。このエントロピー
値が、ゼータ関数の開集合値をその同値類として、ゼータ関数と量子群を加群とする
ための仕組みを形づくるラプラス方程式それ自体が、原子核における中間子論を
形成していると言える。

 一般に言われているゼータ関数は、実軸 $Re={1 \over 2} $に自明な零点が存在
していることが、物理学を持ち出すと、このReが均配力0となり、
相対性理論による慣性質量と重力質量とが同値といえる、慣性系が一定と同じことを
言っている。$Re = {1 \over 2} $は、非対称性理論を言っている。この値の領域
は、遷移元素の光量子仮説から、ゼータ関数と量子群から、原子の固有エントロピー値
を言っている。開集合から、リーマン予想と深リーマン予想を形成するサーストン
多様体をこの開集合の値から、エネルギー安定性が言える。
 
 まとめると、中間子論は、原子核における素粒子同士を結合させるエネルギー
安定性のことで、この粒子は、ゼータ関数と量子群の演算子からラプラス
方程式をつくり出す、大域的微分方程式の役割をする、初期関数を微分演算子
とする、記号論が、この中間子論を形成している。
湯川秀樹教授が、発見した中間子論は、数学の世界では、演算子における記号論
と同値と言えるだろう。

 ラプラス方程式から導かれるヒッグス場の密度エネルギー値は、大域的微分方程式
から求まる、加速度による非線形方程式と、ヴェイユ予想と同値から、
２次曲線による因数定理から、モジュール予想も求まり、その上に、
非慣性系から、無数の接空間からの切片による、D-braneの２次元膜が、
この重力質量に行きつく。

 これらのゼータ関数から導かれる方程式群を調べると、AdS5次元多様体は、
ホログラム空間と言えて、Pholographic spaceとして、宇宙と異次元空間の
アカシックレコードが、心理学だけでは、説明できないことが、数学と物理学
を合わせて、統一場理論を考えていくと、命の樹と命の水が、宇宙の中心にあること
が、分かる。タイムマシーンは、密度エネルギーのコントロールを使うと、
宇宙の周りが過去で、中心が現在であることと、前節の説明から、
未来が宇宙の中心にあるから、命の樹と命の水は、アカシックレコードの源泉
といえるのもわかる。

\end{document}
